\chapter{结论与展望}

\section{结论}

本文将基于人机交互强化的固定翼飞行器着陆控制系统作为研究对象，围绕系统组成、需求分析、动力学模型和仿真平台建立、多工况仿真结果分析以及人机交互强化策略设计等方面展开了研究，主要结论如下：

（1）固定翼飞行器着陆阶段受速度、下沉率、姿态角、侧偏量以及地形环境等多种因素共同影响，是飞行全过程中控制要求较高、风险较集中的阶段。通过建立着陆控制动力学模型和仿真平台，可以较好地反映飞行器在进近、拉平和接地区域的状态变化规律。基准工况校核结果表明，所建模型具有较好的可运行性和合理性，能够满足后续工况分析需要。

（2）多工况仿真结果显示，不同地形条件下飞行器着陆末段的响应呈现显著差别。平坦地形工况整体具有最高的稳定性，可作为参考基线；谷地地形工况表现较好；起伏地形和山脊地形会增加姿态修正频率；阶跃地形对近地判断和接地控制的影响最为突出，属于最为典型的不利工况。地形环境改变对固定翼飞行器着陆控制品质具有重要作用。

（3）本文形成了一种人机交互强化思路，将风险识别、状态提示、操纵辅助和约束保护整合在一起。该策略可在一定程度上提高操纵者对复杂工况的感知和应对能力，对改善着陆安全性和稳定性具有积极作用。

\section{展望}

由于研究条件和论文篇幅限制，本文存在以下不足之处，有待后续深入研究：

（1）本文主要运用仿真分析方法，对真实飞行器复杂的气动特性和实际扰动环境的考量不够全面。后续可结合更高保真度模型进一步优化研究。

（2）本文针对起落架接地缓冲过程主要做了简化处理，后续可结合接地冲击和缓冲系统动力学展开更细致的分析。

（3）本文提出的人机交互强化策略主要在仿真平台上得到验证，后续可结合半实物仿真或飞行试验进一步检验其有效性。

当前，针对固定翼飞行器着陆控制系统的人机交互强化研究领域存在可观的拓展空间。在后续发展中，可以将更复杂的工况条件、更高精度的模型、更完善的交互设计方法相结合，以提高固定翼飞行器在复杂环境中的着陆安全程度，提升其在工程应用方面的价值。
